\section{Free-Agent Threads}
\label{sec:free_agent}
\index{free-agent!free-agent threads}

A logical thread pool exists for every contention group. Prior to OpenMP
Specification 6.0, only assigned
threads from this pool that are part of a team associated with a \kcode{parallel} region are 
available to execute tasks that bind to that \kcode{parallel} region. OpenMP
Specification 6.0 refers to
such threads as \emph{structured threads}. 
OpenMP Specification 6.0 allows threads from the pool that are not assigned to any team to also execute
certain tasks, and refers to such executing threads as \emph{free-agent threads}. 
A task is eligible to be executed by a free-agent thread, and not just by a structured thread 
from the current team, if it is generated by a construct that specifies the \kcode{threadset} 
clause with the \kcode{omp_pool} argument. If the \kcode{threadset} clause is not specified on a construct 
on which it may appear, then the effect is as if the \kcode{threadset} clause was specified 
with \kcode{omp_team} as its argument. This will cause the tasks to be executed only by the
threads of the current team.
Note that there is no means to actually force a given task to be executed by a free-agent thread.

The size of the thread pool (i.e., the thread limit) can be controlled by setting the 
\kcode{OMP_THREAD_LIMIT} environment variable, or by using the \kcode{thread_limit} clause on 
certain constructs. OpenMP Specification 6.0 adds two additional ICVs for setting a limit on the number of active
structured threads versus free-agent threads in a given thread pool. 
With the use of the \kcode{OMP_THREADS_RESERVE} environment variable, these limits can be 
modified by exclusively reserving some number of threads for execution as structured threads or 
for execution as free-agent threads. By default, the behavior is as if 1 thread is exclusively 
reserved for structured thread execution and 0 threads are exclusively reserved for free-agent thread 
execution. This means that the limit on the number of structured threads is equal to the thread 
limit, while the limit on the number of free-agent threads is equal to the thread limit 
minus 1 (at least 1 thread, the ``initial thread" from any thread pool, must be reserved for execution 
as a structured thread). 

In the following example, the \kcode{OMP_THREAD_LIMIT} environment variable is set to 8. Using the \kcode{OMP_THREADS_RESERVE}
environment variable structured thread reservation is set to 5 and free-agent thread reservation is set to 2. 
Though the actual number of structured and unassigned threads are ultimately decided by the implementation, 
the reservation for structured threads has a limiting effect on the number of free-agent threads and vice versa. 
Hence, the assigned threads for any implicit parallel regions cannot exceed 6 (8 minus 2) and the 
free-agent threads cannot exceed 3 (8 minus 5) threads.

\begin{boxedcode}
export OMP_THREAD_LIMIT=8
export OMP_THREADS_RESERVE="structured(5),free_agent(2)"
\end{boxedcode}

Even if no free-agent threads are reserved, an implementation may still provide free-agent threads. 
In the following example, the first task calls the \ucode{do_background_io} procedure, which may
utilize the assigned thread (thread 0) as well as unassigned threads to complete data transfers, 
while the assigned threads execute the parallel region. The \kcode{taskwait} will force the completion 
of all child-tasks including the tasks executed by free-agent threads.

\cexample[6.0]{free_agent}{1}
\ffreeexample[6.0]{free_agent}{1}

In the following example OpenMP tasks are used to calculate the Fibonacci value using both structured 
and free-agent threads. If the task is executed by a free-agent thread then it is recorded by the 
\ucode{count} variable by checking if a free-agent thread is executing the enclosing task region using 
the \kcode{omp_is_free_agent} routine. The thread number of unassigned threads is an implementation-defined 
constant set to \kcode{omp_unassigned_thread} which is a constant less than -1. Even when the use of 
a free-agent thread is allowed via the \kcode{omp_pool} argument, an implementation is not required to use one. 

\cexample[6.0]{free_agent}{2}

\ffreeexample[6.0]{free_agent}{2}

In the next example the \kcode{threadset} clause is used on the \kcode{taskloop} construct. The effect of 
a \kcode{threadset} clause on a \kcode{taskloop} construct is as if it is applied to all the generated tasks. 
One of the two assigned threads and free-agent threads may execute the tasks generated by the \kcode{taskloop} construct. 

\cexample[6.0]{free_agent}{3}

\ffreeexample[6.0]{free_agent}{3}
